\documentclass[11pt]{article}
% ===========================================================================
%  Shared preamble for the AP Calculus AB exam-preparation worksheets.
%  Input right after \documentclass. Files beginning with "_" are not compiled
%  on their own. Navy and gold, to match the course's lesson pages.
% ===========================================================================
\usepackage[letterpaper,top=2.3cm,bottom=2.2cm,left=1.7cm,right=1.7cm,headheight=28pt,footskip=30pt]{geometry}
\usepackage{amsmath,amssymb}
\usepackage[T1]{fontenc}
\usepackage[utf8]{inputenc}
\usepackage{lmodern}
\usepackage{xcolor}
\usepackage[most]{tcolorbox}
\usepackage{enumitem}
\usepackage{fancyhdr}
\usepackage{array}
\usepackage{tikz}
\usepackage{pgfplots}
\pgfplotsset{compat=1.18}

\definecolor{navy}{HTML}{161E54}
\definecolor{ink}{HTML}{27357E}
\definecolor{gold}{HTML}{B98A00}
\definecolor{goldbg}{HTML}{FFFDF3}
\definecolor{inkbg}{HTML}{F6F7FD}
\definecolor{rule}{HTML}{C9CFF0}
\definecolor{okgreen}{HTML}{1B7A44}
\definecolor{okbg}{HTML}{F1F8F3}

\pagestyle{fancy}
\fancyhf{}
\fancyhead[L]{\footnotesize Name: \underline{\hspace{4cm}}}
\fancyhead[C]{\textbf{\color{navy}Integration Academy}\\[-2pt]{\scriptsize Calculus AB \textperiodcentered{} AP\textsuperscript{\textregistered} Exam Preparation}}
\fancyfoot[L]{\tiny\color{black!55}\parbox[t]{0.78\textwidth}{AP\textsuperscript{\textregistered} is a registered trademark of the College Board, which was not involved in the production of, and does not endorse, this material. \textcopyright{} Integration Academy.}}
\fancyfoot[R]{\scriptsize\thepage}
\renewcommand{\headrulewidth}{1.2pt}
\renewcommand{\headrule}{\hbox to\headwidth{\color{gold}\leaders\hrule height \headrulewidth\hfill}}
\renewcommand{\footrulewidth}{0pt}
\setlength{\parindent}{0pt}
\setlength{\parskip}{4pt}

% Title block: \wstitle{code}{title}{unit line}
\newcommand{\wstitle}[3]{\begin{center}{\small\bfseries\color{gold}WORKSHEET #1}\\[3pt]{\LARGE\bfseries\color{navy}#2}\\[3pt]{\small\color{black!60}#3}\end{center}\vspace{2pt}}
% Section heading
\newcommand{\sech}[1]{\par\vspace{10pt}{\large\bfseries\color{navy}#1}\par\nobreak\vspace{2pt}{\color{rule}\hrule height0.8pt}\vspace{7pt}}
% A level tag printed beside a group of questions
\newcommand{\level}[1]{\par\vspace{6pt}{\footnotesize\bfseries\color{gold}\MakeUppercase{#1}}\par\nobreak\vspace{2pt}}

\newtcolorbox{ideabox}{colframe=ink,colback=inkbg,coltitle=white,colbacktitle=ink,fonttitle=\bfseries\small,title={Key ideas},arc=2pt,boxrule=0.8pt,left=8pt,right=8pt}
\newtcolorbox{example}[1]{colframe=rule,colback=white,coltitle=navy,colbacktitle=inkbg,fonttitle=\bfseries\small,title={#1},arc=2pt,boxrule=0.8pt,borderline west={3pt}{0pt}{ink},left=10pt,right=8pt}
\newtcolorbox{frbox}[1]{colframe=gold,colback=goldbg,coltitle=white,colbacktitle=gold,fonttitle=\bfseries\small,title={#1},arc=2pt,boxrule=0.8pt,breakable,left=8pt,right=8pt}
\newtcolorbox{solbox}{colframe=okgreen,colback=okbg,coltitle=white,colbacktitle=okgreen,fonttitle=\bfseries\small,title={Worked solutions},arc=2pt,boxrule=0.8pt,breakable,left=8pt,right=8pt}
\newtcolorbox{scorebox}[1]{colframe=navy,colback=white,coltitle=white,colbacktitle=navy,fonttitle=\bfseries\small,title={#1},arc=2pt,boxrule=0.8pt,breakable,left=8pt,right=8pt}

\newcommand{\soln}{\par\smallskip\textbf{\color{ink}Solution.}\ }
\newcommand{\result}[1]{\par\smallskip\textbf{#1}}
% blank working space of a given height
\newcommand{\work}[1]{\par\vspace{#1}}
% a point value in the margin of a scoring guide
\newcommand{\pts}[1]{\hfill{\footnotesize\bfseries\color{gold}[#1]}}

% Coordinate grid for graph questions: \begin{apgraph}{xmin}{xmax}{ymin}{ymax} ... \end{apgraph}
\newenvironment{apgraph}[4]{\begin{center}\begin{tikzpicture}\begin{axis}[
  axis lines=middle, xmin=#1, xmax=#2, ymin=#3, ymax=#4,
  grid=both, grid style={black!12}, xtick distance=1, ytick distance=1,
  ticklabel style={font=\scriptsize}, xlabel={$x$}, ylabel={$y$},
  width=10cm, height=7.2cm, samples=80, clip=false,
  every axis plot/.append style={line width=1.1pt, color=ink}]}{\end{axis}\end{tikzpicture}\end{center}}
\newcommand{\closeddot}[2]{\addplot[only marks, mark=*, mark size=2.2pt, color=ink] coordinates {(#1,#2)};}
\newcommand{\opendot}[2]{\addplot[only marks, mark=*, mark size=2.2pt, color=ink, mark options={fill=white}] coordinates {(#1,#2)};}

% Multiple-choice options: four across, or stacked when the options are long.
\newcommand{\opts}[4]{\par\vspace{3pt}\begin{tabular}{@{}p{0.235\linewidth}p{0.235\linewidth}p{0.235\linewidth}p{0.235\linewidth}@{}}(A)\ #1 & (B)\ #2 & (C)\ #3 & (D)\ #4\end{tabular}\par\vspace{9pt}}
\newcommand{\optsv}[4]{\par\vspace{3pt}\begin{tabular}{@{}l@{\ \ }p{0.9\linewidth}@{}}(A) & #1\\(B) & #2\\(C) & #3\\(D) & #4\end{tabular}\par\vspace{9pt}}

\fancyhead[R]{\footnotesize\color{black!60}Worksheet 8.6}
\begin{document}

\wstitle{8.6}{Regions with More Than One Piece}{Unit 8: Applications of Integration}

\begin{ideabox}
\begin{itemize}[leftmargin=*,itemsep=2pt,topsep=2pt]
  \item If the graphs cross inside the interval, the top graph changes. Split the integral at each crossing.
  \item In one formula: area $=\displaystyle\int_{a}^{b}|f(x)-g(x)|\,dx$. On a calculator, enter the absolute value directly.
  \item $\displaystyle\int_{a}^{b}\big(f(x)-g(x)\big)\,dx$ without the absolute value gives a \textbf{net} amount, in which pieces cancel.
  \item A region with a corner in its boundary also needs two integrals, one for each boundary curve.
\end{itemize}
\end{ideabox}

\sech{Worked examples}
\begin{example}{Example 1 \textperiodcentered{} A cubic and a line}
Find the total area enclosed by $y=x^3$ and $y=4x$.\soln They meet at $x=-2,0,2$. By symmetry, $2\displaystyle\int_{0}^{2}(4x-x^3)\,dx$.\result{$2(8-4)=8$.}
\end{example}

\begin{example}{Example 2 \textperiodcentered{} Net against total}
Compare $\displaystyle\int_{0}^{\pi}\cos x\,dx$ with the area between $y=\cos x$ and the $x$-axis for $0\le x\le\pi$.\soln The integral is 0. The area is $\int_{0}^{\pi/2}\cos x\,dx+\left|\int_{\pi/2}^{\pi}\cos x\,dx\right|$.\result{The integral is 0; the area is 2.}
\end{example}

\begin{example}{Example 3 \textperiodcentered{} Splitting at a zero}
Find the area between $y=x^2-4$ and the $x$-axis for $0\le x\le3$.\soln $\displaystyle\int_{0}^{2}(4-x^2)\,dx=\dfrac{16}{3}$ and $\displaystyle\int_{2}^{3}(x^2-4)\,dx=\dfrac73$.\result{$\dfrac{23}{3}$.}
\end{example}

\newpage
\sech{Practice}

\level{Core}
\begin{enumerate}[leftmargin=*,itemsep=2pt]
  \item Find the area between $y=x-1$ and the $x$-axis for $0\le x\le4$.\work{2.3cm}
  \item At what value of $x$ in $(0,3)$ do the graphs of $y=x^2$ and $y=2x$ cross?\work{2.3cm}
  \item Find the area between $y=x^2$ and $y=2x$ for $0\le x\le3$.\work{2.3cm}
  \item Evaluate $\displaystyle\int_{0}^{3}(x^2-2x)\,dx$, and explain the result using Question 3.\work{2.3cm}
\end{enumerate}

\level{Exam level}
\begin{enumerate}[leftmargin=*,itemsep=2pt,start=5]
  \item Find the total area between $y=x^3-x$ and the $x$-axis for $-1\le x\le1$.\work{2.8cm}
  \item Find the area between $y=\sin x$ and $y=\dfrac12$ for $0\le x\le\dfrac\pi2$.\work{2.8cm}
  \item Find the area of the region in the first quadrant bounded by $y=x^2$, $y=2-x$ and the $x$-axis.\work{2.8cm}
  \item \emph{Calculator.} Find the area between $y=\cos x$ and $y=\dfrac x2$ for $0\le x\le3$.\work{2.8cm}
\end{enumerate}

\begin{frbox}{Question 9 \textperiodcentered{} Free-response practice}
Let $f(x)=x^3-5x^2+4x$.
\begin{enumerate}[label=(\alph*),leftmargin=*,itemsep=2pt]
  \item Evaluate $\displaystyle\int_{0}^{4}f(x)\,dx$.\work{2.6cm}
  \item Find the total area between the graph of $f$ and the $x$-axis for $0\le x\le4$.\work{2.6cm}
  \item Explain why the answers to parts (a) and (b) differ, and write a single integral that gives the area in part (b).\work{2.6cm}
\end{enumerate}
\end{frbox}

\newpage
\begin{solbox}
\begin{enumerate}[leftmargin=*,itemsep=4pt]
  \item $\dfrac12+\dfrac92=5$.
  \item $x=2$.
  \item $\displaystyle\int_{0}^{2}(2x-x^2)\,dx+\displaystyle\int_{2}^{3}(x^2-2x)\,dx=\dfrac43+\dfrac43=\dfrac83$.
  \item $0$: the two pieces have equal area, and they cancel.
  \item $2\left(\dfrac14\right)=\dfrac12$.
  \item They cross at $x=\dfrac\pi6$: $\left(\dfrac{\pi}{12}-1+\dfrac{\sqrt3}{2}\right)+\left(\dfrac{\sqrt3}{2}-\dfrac\pi6\right)=\sqrt3-1-\dfrac{\pi}{12}$.
  \item The curves meet at $x=1$: $\displaystyle\int_{0}^{1}x^2\,dx+\displaystyle\int_{1}^{2}(2-x)\,dx=\dfrac13+\dfrac12=\dfrac56$.
  \item $\displaystyle\int_{0}^{3}\left|\cos x-\tfrac x2\right|\,dx\approx3.293$ (the graphs cross at $x\approx1.030$).
  \item \begin{enumerate}[label=(\alph*),leftmargin=*,itemsep=3pt]
    \item $\Big[\tfrac{x^4}{4}-\tfrac{5x^3}{3}+2x^2\Big]_0^4=-\dfrac{32}{3}$.
    \item $f(x)=x(x-1)(x-4)$: $\displaystyle\int_{0}^{1}f(x)\,dx=\dfrac{7}{12}$ and $\displaystyle\int_{1}^{4}f(x)\,dx=-\dfrac{45}{4}$. Area $=\dfrac{7}{12}+\dfrac{45}{4}=\dfrac{71}{6}$.
    \item Part (a) counts the piece below the axis as negative, so the pieces partly cancel. Area $=\displaystyle\int_{0}^{4}|f(x)|\,dx$.
  \end{enumerate}
\end{enumerate}
\end{solbox}

\end{document}
